\section{Conclusions and Discussion}
This paper introduced a graph-based feature extraction methodology for Constraint Programming instances, with a specific focus on cut-based representations (\texttt{WLc}). By integrating this process directly into the FlatZinc ecosystem, we bridged the gap between structural graph representations and classical machine learning paradigms.

Our findings lead to several key conclusions regarding the efficacy of these features:
(i) \textbf{Robustness Across Paradigms:} The \texttt{WL} and \texttt{WLc} features demonstrate high versatility, providing strong predictive signals across geometric (SVM), logical (RF), and deep (MLP) models. This suggests that the structural information captured is fundamental and not dependent on a specific learning architecture. (ii) \textbf{Stability and Performance Reliability:} The proposed features exhibit high stability across diverse tasks. Even in scenarios where statistical significance over the baseline was not reached, the mean and median performance scores remained consistently higher. Notably, the \texttt{WLc} features do not suffer from the underperformance observed in traditional statistical baselines during specific tasks, such as parallel improvement prediction with Random Forests. (iii) \textbf{State-of-the-Art Accuracy:} For the task of predictive accuracy maximization this methodology establishes a new performance benchmark. The cut-based features consistently outperformed \texttt{fzn2feat} with high statistical significance, identifying the optimal solver with greater reliability than existing statistical methods.

Future work will focus on optimizing the computational efficiency of the extraction pipeline and exploring the integration of these features into dynamic, online algorithm selection frameworks. Furthermore we will explore the effect of combining our feature sets with classical feature sets (like \texttt{fzn2feat}) to evaluate if combining different types of signals can improve the predictive performance. We will also focus on generalise our features to different tasks not strictly related to algorithm selection but still relevant in the CP community.   